\documentclass[../../../include/open-logic-section]{subfiles}

\begin{document}

\olfileid{sfr}{siz}{nen-alt}

\olsection{不可数的集合}

\begin{editorial}
本节使用 \olref[enm-alt]{sec} 中的定义（即要求与~$\Nat$ 建立双射，而非从 $\PosInt$ 到其上的满射）来证明 $\Bin^\omega$ 和 $\Pow{\Nat}$ 的不可数性。
\end{editorial}

% Part:sets-functions-relations
% Chapter: size-of-sets
% Section: non-enumerability-alt

\begin{explain}
自然数集 $\Nat$ 是无穷的。它显然也是可数的。
但一个引人注目的事实是，存在\emph{不可数的}集合，即不可数的集合
（参见\olref[sfr][siz][enm-alt]{defn:enumerable}）。

这可能会令人惊讶。毕竟，说 $A$ 不可数，就是说\emph{不存在}双射 $f
\colon \Nat \to A$；也就是说，没有任何将 $\Nat$ 的无穷多个元素映射到 $A$ 的函数能穷尽 $A$。
因此，如果 $A$ 是不可数的，那么 $A$ 的元素就比自然数“更多”。

要证明一个集合不可数，就必须证明不可能存在适当的双射。
最好的方法是证明，任何枚举 $A$ 的元素的尝试都必定至少遗漏一个元素；这表明没有任何函数 $f\colon \Nat
\to A$ 是满射。而确立这一点的一个通用策略是使用 Cantor（康托尔）的\emph{对角线法}。
给定 $A$ 的元素的一个列表，比如 $x_1$, $x_2$, \dots，我们构造 $A$ 的另一个元素，根据其构造方式，它不可能出现在该列表中。

不过，通过例子来理解这一切是最好的方式。所以，我们的第一个例子是所有由 $0$ 和 $1$ 组成的无穷序列的集合 $\Bin^\omega$。
（‘$\Bin$’代表 binary（二进制），我们可以简单地将它看作二元集
$\{0,1\}$。）\oliflabeldef{sfr:card-arithmetic:card-opps:sec}{\footnote{更准确地说，我们应当规定 $\Bin^\omega$ 是所有由 $0$ 和 $1$ 组成的 $\omega$-序列的集合，即集合
$\funfromto{\omega}{\{0,1\}}$。但其确切含义要到
\olref[sfr][card-arithmetic][card-opps]{sec} 中才会清晰。}{}不过，目前这种稍欠严格的表述应该不会引起任何混淆。}
\end{explain}

\begin{thm}
\ollabel{thm:nonenum-bin-omega}
$\Bin^\omega$~是不可数的。
\end{thm}

\begin{proof}
考虑 $\Bin^\omega$ 的某个子集的任意一种枚举。于是我们有一个列表 $s_{0}$, $s_{1}$, $s_{2}$, \dots{}，其中每个 $s_n$ 都是由 $0$ 和~$1$ 组成的无穷序列。令 $s_n(m)$ 为该列表中第 $n$ 个字符串的第 $m$ 位数字。这样我们就可以将列表视为一个阵列，其中 $s_n(m)$ 位于第 $n$ 行第 $m$ 列：
\[
\begin{array}{c|c|c|c|c|c}
& 0 & 1 & 2 & 3 & \dots \\\hline
0 & \mathbf{s_{0}(0)} & s_{0}(1) & s_{0}(2) & s_0(3) & \dots \\\hline
1 & s_{1}(0)& \mathbf{s_{1}(1)} & s_1(2) & s_1(3) & \dots \\\hline
2 & s_{2}(0)& s_{2}(1) & \mathbf{s_2(2)} & s_2(3) & \dots \\\hline
3 & s_{3}(0)& s_{3}(1) & s_3(2) & \mathbf{s_3(3)} & \dots \\\hline
\vdots & \vdots & \vdots & \vdots & \vdots & \mathbf{\ddots}
\end{array}
\]
现在，我们要构造一个由 $0$ 和 $1$ 组成的无穷序列 $d$，它不在此列表中。我们将通过指定其每一项来完成构造，即对所有 $n \in \Nat$ 指定 $d(n)$。直观地说，我们沿着上述阵列的对角线向下读取（因此得名“对角线法”），然后将 $1$ 变成 $0$，$0$ 变成 $1$ 来构造 $d$。更抽象地说，我们根据对角线上的第 $n$ 个元素 $s_n(n)$ 是 $1$ 还是 $0$，来定义 $d(n)$ 为 $0$ 或 $1$，即：
\[
d(n) =
\begin{cases}
1 & \text{当$s_{n}(n) = 0$}\\
0 & \text{当$s_{n}(n) = 1$}
\end{cases}
\]
显然 $d \in \Bin^\omega$，因为它是由 $0$ 和 $1$ 组成的无穷序列。但我们构造 $d$ 的方式使得对于任意 $n \in \Nat$，都有$d(n) \neq s_n(n)$。也就是说，$d$ 与 $s_n$ 在第 $n$ 项上不同。因此，对任意 $n\in \Nat$，$d \neq s_n$。所以 $d$ 不可能在列表 $s_0$, $s_1$, $s_2$, \dots 中。

我们已经证明，对于 $\Bin^\omega$ 的任意子集的任意枚举，它都会遗漏 $\Bin^\omega$ 中的某个元素。因此，集合 $\Bin^\omega$ 不存在枚举，即 $\Bin^\omega$ 是不可数的。
\end{proof}

\begin{explain}
这种证明方法之所以被称为“对角线化”，是因为它利用阵列的对角线来定义~$d$。然而，对角线化并不一定需要出现阵列。事实上，即使没有阵列和实际的“对角线”参与，我们也可以使用类似的思想来证明某个集合是不可数的。下面的结果就说明了这一点。
\end{explain}

\begin{thm}
\ollabel{thm:nonenum-pownat}
$\Pow{\Nat}$ 不是可数的。
\end{thm}

\begin{proof}
我们采用同样的方式进行证明，即证明任何 $\Nat$ 的子集列表都会遗漏某个 $\Nat$ 的子集。假设我们有一个 $\Nat$ 的子集列表 $N_0, N_1, N_2, \ldots$。我们如下定义集合 $D$：$n \in D$ 当且仅当 $n \notin N_{n}$：
\[
D = \Setabs{n \in \Nat}{n \notin N_n}
\]
显然 $D\subseteq \Nat$。但 $D$ 不可能在此列表中。毕竟，根据构造，$n \in D$ 当且仅当 $n\notin N_n$，因此对于任意 $n \in \Nat$，$D \neq N_n$。
\end{proof}

\begin{explain}
前面的证明并未明确提及对角线。不过，如果你这样想象，仍可以认为它涉及了对角线：想象将集合 $N_0$, $N_1$, \dots 写成一个阵列，其中我们在第 $n$ 行写下 $N_n$，方式是当 $m \in N_n$ 时在第 $m$ 列写下 $m$。例如，假设该列表上的前四个集合是 $\{0,1,2,\dots\}$, $\{1, 3, 5, \dots\}$, $\{0,1,4\}$ 和 $\{2,3,4,\dots\}$；那么我们的阵列将以
\[
\begin{array}{r@{}rrrrrrr}
  N_0 = \{ & \mathbf{0}, & 1, & 2, & & & & \dots\}\\
  N_1 = \{ &  & \mathbf{1}, &  & 3, &  & 5, & \dots\}\\
  N_2 = \{ & 0, & 1, &  &  & 4\phantom{,} &  & \}\\
  N_3 = \{ &  &  & 2, & \mathbf{3}, & 4, & & \dots\}\\
  &\vdots & & & & & & \ddots\phantom{\}}
  \end{array}
\]
开头。那么，$D$ 就是沿对角线向下遍历得到的集合，其中 $n \in D$ 当且仅当 $n$ \emph{不}在对角线上。因此，在上述例子中，我们会排除 $0$ 和 $1$，纳入 $2$，排除 $3$，等等。
\end{explain}

\begin{prob}
通过显式的对角线论证，证明所有函数 $f \colon \Nat \to \Nat$ 的集合是不可数的。也就是说，证明如果 $f_1$, $f_2$, \dots 是一个函数列表，且每个 $f_i\colon \Nat \to \Nat$，那么存在某个 $g \colon \Nat \to \Nat$ 不在此列表中。
\end{prob}

\end{document}